\scanpage{103}

\uline{Corollaire} $N(u)=N(|u|)$\,; si $u=h+ik$ $N(h),N(k)\leq N(u)\leq N(h)+N(k)$, alors $u\in L^N(E)\deleted{\unsure{F}}$ \deleted{$\Longleftrightarrow$}$\Longleftrightarrow h,k\in L^N(E)\deleted{\unsure{(F)}}$. \deleted{\unsure{et}}

\margin{Analogue pour $L_N(E,F)$}

\uline{Prop} $L^N(E,F)\subset L_0(E,F)$ si $N$ non équivalente à $N_\infty$\unsure{;} (sinon $L^N(E,F)=L(E,F)$). En effet, si $u\in L^N(E,F)$, \deleted{prouvons que}\inserted{alors} $|u|\in L^N(E)$, prouvons $|u|$ compact (d'où $u$ compact). Sinon, il existerait un op. hermitien $v\in L(E)$ tel que $v|u|$ soit un projecteur hermitien infini, or on aurait $v|u|\in L^N(E)$. C'est impossible, car on \deleted{\unsure{aura|aurait} $(\rho_i)\in\ell^N$} aurait $\sum|x_i|\deleted{\unsure{1}}\leq M$ pour $N^\circ((x_i))\leq 1$, donc $N^\circ$ équivalente à $N_1$, ce qui est absurde.

\uline{Corollaire} $L^{N_p}(E,F)=L^p(E,F)$ pour $1\leq p\leq+\infty$\,; \deleted{si $p<+\infty$, est}\inserted{si $p<+\infty$, est} même identique à $L_{N_p}(E,F)$. (car $\ell_N=\ell^N\unsure{=|\Longrightarrow}L_N(E,F)=L^N(E,F)$\unsure{)}

\uline{Prop} $u\in L^N(E,F)$, $v\in L^{N^\circ}(F,G)\Rightarrow vu\in L^1(E,G)$

D'où accouplement Trace entre $L^N(E,F)$ et \deleted{\unsure{sur $C$|$\to C$}} $L^{N^\circ}(F,E)$. Par restriction, \deleted{$L^{N^\circ}(F,E)$ est un sous-espace du dual} en principe, on a

\uline{Prop} le dual de $L_N(E,F)$ est $L^N(E,F)$

\uline{Corollaire 1} $L^N(E,F)$ et $L_N(E,F)$ sont complets

\uline{Corollaire 2} $L_N(E,F)$ réflexif $\Leftarrow$ \uline{$N$ réflexif}. \margin{\begin{tikzpicture}\draw[red,line width=0.8pt] (0,0)--(0.01,0.65);\end{tikzpicture}} Réciproque vraie si $E$ et $F$ de dimension infinie. \uline{En tous cas}, $L_N(E,F)$ réfl. $\Longleftrightarrow L^N(E,F)$ réflexif $\Longleftrightarrow\ldots$

\uline{Corollaire 3} Si $1<p<+\infty$, alors $L^p(E,F)$ est un espace réflexif ayant pour dual $L^{p'}(F,E)$

\uline{Prop} Si $N,N',N''$ sont des normes de Schatten telles que $N(xy)\leq N'(x)N'(y)$, alors $u\in L^{N'}(E,F)$, $v\in L^{N''}(F,G)$ implique $vu\in L^N(E,F)$, \inserted{et $N(vu)\leq N'(u)N''(v)$ \uline{Corollaire}} Si \uline{un des \deleted{\unsure{soit}} $N',N''$ est mis en} bas, on peut faire de $=$ avec l'autre. \margin{\begin{tikzpicture}\draw[red,line width=0.8pt] (0,0)--(0.04,1.15);\end{tikzpicture}}

(\deleted{par transposition, on se ramène \unsure{au deuxième} \illegible{2}, en} \deleted{prenant $N'''=N^\circ$,} Par le deuxième th. de convexité si $v$ et $u$ sont \deleted{\unsure{là par contraction}, on} \deleted{à prouver que si $N_1(xyz)\leq N'(x)N''(y)N'''(z)$, alors} \unsure{cpts.|cpcts.} On s'y ramène facilement, en notant que \deleted{on a $N_1(wvu)\leq N'(u)N''(v)N'''(w)$.}

\[
N(u)=\lim N(\underline{pu}\deleted{\unsure{p|q}})=\lim_p N(pup)=\lim_p N(\unsure{up|u_p})
\]

$p$ projecteur orth. de rang fini, \uline{variant} en filtrant croissant. (\unsure{Car|Donc} $N(pvup)=N((pv)(up))\leq N''(pv)N^{\unsure{''|'}}(up)\leq N''(v)N^{\unsure{\circ|'}}(u)$). Enfin, rappelons

\uline{Prop} $N((\lambda_i u))\leq N(u)$

\uline{Prop} Caractérisation des espaces de Schatten.

\margin{104}
